\documentclass[a4paper,12pt]{article}

\usepackage[T2A]{fontenc}
\usepackage[utf8]{inputenc}
\usepackage[russian]{babel}
\usepackage{amsmath}
\usepackage{PTSans}
\renewcommand{\familydefault}{\sfdefault}
\usepackage{icomma}
\usepackage[a4paper,margin=18mm,headheight=15pt]{geometry}
\usepackage{microtype}
\usepackage{tikz}
\usetikzlibrary{calc,arrows.meta,positioning,shapes.geometric}
\usepackage{tabularx,booktabs}
\usepackage{enumitem}
\usepackage{parskip}
\usepackage{fancyhdr}
\usepackage{setspace}
\usepackage{xcolor}

\definecolor{accent}{HTML}{008080}
\definecolor{darkaccent}{HTML}{004D4D}
\definecolor{textgray}{HTML}{2F3437}
\definecolor{softgray}{HTML}{6B7280}
\definecolor{linegray}{HTML}{D7DEDE}

\setlength{\parindent}{0pt}
\onehalfspacing
\color{textgray}
\renewcommand{\arraystretch}{1.25}

\setlist[itemize]{
    leftmargin=1.45em,
    itemsep=0.18em,
    topsep=0.28em
}
\setlist[enumerate]{
    leftmargin=1.8em,
    itemsep=0.45em,
    topsep=0.3em
}

\pagestyle{fancy}
\fancyhf{}
\fancyhead[L]{\small Екатерина Гнедкова}
\fancyhead[R]{\small ЕГЭ · текстовые задачи}
\fancyfoot[C]{\small\thepage}
\renewcommand{\headrulewidth}{0.3pt}
\renewcommand{\footrulewidth}{0pt}

\newcommand{\lessondate}{25 сентября 2026~г.}
\newcommand{\sect}[1]{%
    \vspace{0.35em}%
    {\Large\bfseries\color{darkaccent}#1}\par
    \vspace{0.22em}%
    \color{linegray}\hrule
    \color{textgray}%
    \vspace{0.55em}%
}
\newcommand{\subsect}[1]{%
    \vspace{0.3em}%
    {\large\bfseries\color{accent}#1}\par
    \vspace{0.18em}%
}
\newcommand{\important}[1]{\textbf{\color{darkaccent}#1}}
\newcommand{\unit}[1]{\,\text{#1}}

\tikzset{
    flowstart/.style={
        ellipse,
        draw=darkaccent,
        line width=0.9pt,
        minimum width=34mm,
        minimum height=9mm,
        align=center,
        font=\small
    },
    flowproc/.style={
        rounded corners=3pt,
        rectangle,
        draw=accent,
        line width=0.85pt,
        text width=104mm,
        minimum height=11mm,
        align=center,
        font=\small,
        inner sep=4pt
    },
    flowbranch/.style={
        rounded corners=3pt,
        rectangle,
        draw=accent,
        line width=0.85pt,
        text width=56mm,
        minimum height=13mm,
        align=center,
        font=\small,
        inner sep=4pt
    },
    flowdecision/.style={
        diamond,
        aspect=2.15,
        draw=darkaccent,
        line width=0.9pt,
        text width=40mm,
        align=center,
        font=\small,
        inner sep=2pt
    },
    flowarrow/.style={
        -{Stealth[length=3mm,width=2mm]},
        line width=0.9pt,
        draw=textgray
    }
}

\begin{document}

% ============================================================
% ТИТУЛЬНЫЙ ЛИСТ
% ============================================================

\thispagestyle{empty}
\begin{center}

\vspace*{22mm}

{\small\color{softgray}
ПОДГОТОВКА К ЕГЭ ПО ПРОФИЛЬНОЙ МАТЕМАТИКЕ
}\par

\vspace{9mm}

{\Huge\bfseries\color{darkaccent}
Чек-лист
}\par

\vspace{4mm}

{\LARGE\bfseries
Текстовые задачи: движение и производительность
}\par

\vspace{13mm}

{\Large
Екатерина Гнедкова
}\par

\vspace{3mm}

{\large
Формат тренировки: Путь А
}\par

\vspace{15mm}

\begin{minipage}{0.76\textwidth}
\centering
\large
Табличный метод · движение навстречу и вдогонку ·
круговая трасса · движение по реке · средняя скорость ·
протяжённые тела · производительность
\end{minipage}

\vfill

{\normalsize \lessondate}\par

\vspace{4mm}

{\small\color{softgray}
Лёвин Артём Александрович эксклюзивно для Екатерины Гнедковой
}\par

\vspace{16mm}

\begin{tikzpicture}[remember picture,overlay]
    \draw[
        accent!55,
        line width=1.2pt
    ]
    ($(current page.south west)+(18mm,18mm)$)
    .. controls
    ($(current page.south west)+(76mm,30mm)$)
    and
    ($(current page.south east)+(-72mm,7mm)$)
    ..
    ($(current page.south east)+(-18mm,18mm)$);
\end{tikzpicture}

\end{center}

\clearpage

% ============================================================
% ОСНОВНОЙ ЧЕК-ЛИСТ
% ============================================================

\sect{1. Универсальная структура текстовой задачи}

В задачах занятия используется одна и та же идея:
связать \important{величину за единицу времени},
\important{время} и \important{результат}.

\begin{table}[h!]
\centering
\begin{tabularx}{\linewidth}{@{}>{\bfseries}lXXX@{}}
\toprule
Тип задачи & Величина за единицу времени & Время & Результат \\
\midrule
Движение & скорость $v$ & $t$ & путь $S=vt$ \\
Производительность & производительность $p$ & $t$ & работа или объём $A=pt$ \\
\bottomrule
\end{tabularx}
\end{table}

\important{Рабочий принцип:}
сначала обозначьте через $x$ величину, которую требуется найти.
После этого перенесите известные данные в таблицу и восстановите
недостающие клетки по формулам

\[
S=vt,
\qquad
t=\frac{S}{v},
\qquad
v=\frac{S}{t},
\]

или

\[
A=pt,
\qquad
t=\frac{A}{p},
\qquad
p=\frac{A}{t}.
\]

Затем найдите в тексте фразу, которая связывает строки таблицы:
\emph{времена равны}, \emph{одно время больше другого},
\emph{суммарное время равно заданному},
\emph{быстрый прошёл на один круг больше} и т.

\subsect{Единицы измерения}

Все величины в одном уравнении должны быть согласованы.
Если скорость дана в км/ч, удобно переводить минуты в часы,
а метры --- в километры:

\[
40\unit{мин}=\frac{2}{3}\unit{ч},
\qquad
300\unit{м}=0{,}3\unit{км}.
\]

В задачах по математике переход к СИ не обязателен.
Главное --- использовать единую систему единиц внутри расчёта.

\subsect{Ограничения и самопроверка}

Перед решением проверьте: знаменатели не равны нулю; скорость, время и
производительность положительны; при движении против течения собственная
скорость лодки больше скорости течения. После решения сопоставьте найденное
значение с условием и запишите именно ту величину и те единицы, которые требуются в вопросе.

% ============================================================
% БЛОК-СХЕМА
% ============================================================

\clearpage
\thispagestyle{empty}

\begin{center}
{\Large\bfseries\color{darkaccent}
Алгоритм: от текста задачи к уравнению
}\par
\vspace{7mm}

\begin{tikzpicture}[node distance=5.5mm and 13mm]

\node[flowstart] (start)
{Прочитать вопрос задачи};

\node[flowproc,below=of start] (x)
{Обозначить через $x$ искомую величину;
при необходимости ввести вспомогательное обозначение};

\node[flowdecision,below=7mm of x] (type)
{Какой тип задачи?};

\node[flowbranch,below left=10mm and 7mm of type] (move)
{Движение: таблица $v$--$t$--$S$, формула $S=vt$};

\node[flowbranch,below right=10mm and 7mm of type] (work)
{Работа: таблица $p$--$t$--$A$, формула $A=pt$};

\node[flowproc,below=27mm of type] (link)
{Найти связь между строками таблицы:
равенство времён, разность времён, общий путь,
отрыв, один круг, суммарное время};

\node[flowproc,below=of link] (eq)
{Записать ограничения и составить уравнение};

\node[flowproc,below=of eq] (solve)
{Решить уравнение: раскрыть скобки,
сократить дроби, привести к общему знаменателю,
привести подобные};

\node[flowdecision,below=7mm of solve] (check)
{Ответ удовлетворяет условию
и единицам?};

\node[flowstart,below=8mm of check] (finish)
{Записать ответ};

\draw[flowarrow] (start) -- (x);
\draw[flowarrow] (x) -- (type);
\draw[flowarrow] (type.west) -| node[pos=0.25,above,font=\small]{движение} (move.north);
\draw[flowarrow] (type.east) -| node[pos=0.25,above,font=\small]{работа} (work.north);
\draw[flowarrow] (move.south) |- (link.west);
\draw[flowarrow] (work.south) |- (link.east);
\draw[flowarrow] (link) -- (eq);
\draw[flowarrow] (eq) -- (solve);
\draw[flowarrow] (solve) -- (check);
\draw[flowarrow] (check) -- node[right,font=\small]{да} (finish);
\draw[flowarrow] (check.west) -- ++(-18mm,0)
node[midway,above,font=\small]{нет}
|- (eq.west);

\end{tikzpicture}
\end{center}

\clearpage

% ============================================================
% ДВИЖЕНИЕ НАВСТРЕЧУ
% ============================================================

\sect{2. Движение навстречу}

При одновременном движении навстречу друг другу скорость сближения равна

\[
v_{\text{сбл}}=v_1+v_2.
\]

Если начальное расстояние между участниками равно $S$, то время до встречи

\[
t=\frac{S}{v_1+v_2}.
\]

Если один участник выехал раньше, его время в пути больше на величину форы.

\important{Типовой приём:}
если требуется найти расстояние от пункта $A$ до места встречи,
обозначьте это расстояние через $x$.
Тогда второй участник пройдёт расстояние $S-x$.

Например, первый автомобиль движется со скоростью $60\unit{км/ч}$,
второй --- $65\unit{км/ч}$, а первый находился в пути на $1\unit{ч}$ дольше.
Тогда

\[
\frac{x}{60}-\frac{S-x}{65}=1.
\]

Разность записана именно в таком порядке,
поскольку время первого автомобиля больше.

% ============================================================
% ДОГОНКА
% ============================================================

\sect{3. Движение вдогонку}

При движении в одном направлении используется скорость сближения

\[
v_{\text{сбл}}=v_{\text{быстр}}-v_{\text{медл}}.
\]

Если начальный отрыв равен $d$, то

\[
t=\frac{d}{v_{\text{быстр}}-v_{\text{медл}}},
\qquad
v_{\text{быстр}}>v_{\text{медл}}.
\]

Иногда удобно ввести вспомогательную переменную.
Пусть $x$ --- время движения, а $y$ --- скорость медленного пешехода.
Если быстрый идёт на $1{,}5\unit{км/ч}$ быстрее и за это время ликвидирует отрыв $0{,}3\unit{км}$, то

\[
x(y+1{,}5)-xy=0{,}3.
\]

После раскрытия скобок величина $y$ сокращается:

\[
1{,}5x=0{,}3,
\qquad
x=0{,}2\unit{ч}=12\unit{мин}.
\]

\important{Вывод:}
вспомогательная переменная допустима, если после упрощения уравнение
однозначно определяет искомую величину.

% ============================================================
% КРУГОВАЯ ТРАССА
% ============================================================

\sect{4. Круговая трасса}

Если участники стартуют из одной точки и движутся в одном направлении,
быстрый участник обгоняет медленного ровно на один круг,
когда выигрыш в пройденном пути становится равен длине трассы $L$:

\[
\bigl(v_{\text{быстр}}-v_{\text{медл}}\bigr)t=L.
\]

Для обгона на $k$ кругов:

\[
\bigl(v_{\text{быстр}}-v_{\text{медл}}\bigr)t=kL.
\]

Если участники стартуют из противоположных точек окружности,
то расстояние между ними по каждой полуокружности равно

\[
\frac{L}{2}.
\]

При движении навстречу до первой встречи:

\[
(v_1+v_2)t=\frac{L}{2}.
\]

\subsect{Что не перепутать}

При обгоне длина круга появляется как
\important{разность пройденных путей}.
При старте из противоположных точек половина длины круга --- это
\important{начальный промежуток по дуге}.

% ============================================================
% РЕКА
% ============================================================

\sect{5. Движение по реке}

Пусть $v$ --- собственная скорость лодки,
а $u$ --- скорость течения. Тогда

\[
v_{\text{по течению}}=v+u,
\qquad
v_{\text{против течения}}=v-u,
\qquad
v_{\text{плота}}=u.
\]

Для движения против течения необходимо

\[
v>u.
\]

Если лодка проходит расстояние $x$ по течению,
стоит $t_0$ часов и возвращается на то же расстояние против течения,
то полное время рейса равно

\[
\frac{x}{v+u}+t_0+\frac{x}{v-u}.
\]

Если требуется суммарный путь туда и обратно,
после нахождения $x$ необходимо вычислить $2x$.

\important{Неодновременный старт:}
если плот вышел раньше лодки, его временная фора обязательно должна
присутствовать в уравнении. Плот движется только со скоростью течения.

% ============================================================
% СРЕДНЯЯ СКОРОСТЬ
% ============================================================

\sect{6. Средняя скорость}

Средняя скорость вычисляется по определению:

\[
\boxed{
 v_{\text{ср}}=
 \frac{S_{\text{общ}}}{t_{\text{общ}}}
}
\]

Простое среднее арифметическое скоростей допустимо только в специальных случаях,
например при равных промежутках времени.

Если весь путь разделён на три равные по длине части,
а скорости на них равны $v_1$, $v_2$, $v_3$,
обозначим весь путь через $S$.
Тогда каждое расстояние равно $S/3$, а времена равны

\[
\frac{S}{3v_1},
\qquad
\frac{S}{3v_2},
\qquad
\frac{S}{3v_3}.
\]

Поэтому

\[
v_{\text{ср}}
=
\frac{S}{
\dfrac{S}{3v_1}
+
\dfrac{S}{3v_2}
+
\dfrac{S}{3v_3}}
=
\frac{3}{
\dfrac{1}{v_1}
+
\dfrac{1}{v_2}
+
\dfrac{1}{v_3}}.
\]

\important{Контрольный вопрос:}
что одинаково --- расстояния или времена?
От этого зависит способ вычисления средней скорости.

% ============================================================
% ПРОТЯЖЁННЫЕ ТЕЛА
% ============================================================

\sect{7. Движение протяжённых тел}

В задачах с поездами, платформами, мостами и другими объектами,
размерами которых нельзя пренебречь, удобно рассматривать относительное движение.

Если поезд длиной $L_1$ полностью проходит мимо неподвижной платформы длиной $L_2$,
то от момента, когда локомотив поравнялся с началом платформы,
до момента, когда последний вагон миновал её конец,
поезд проходит расстояние

\[
L_1+L_2.
\]

Следовательно,

\[
t=\frac{L_1+L_2}{v}.
\]

Если два протяжённых объекта полностью проходят друг мимо друга,
то в числителе появляется сумма их длин.
При движении навстречу

\[
t=\frac{L_1+L_2}{v_1+v_2},
\]

а при движении в одном направлении

\[
t=\frac{L_1+L_2}{v_{\text{быстр}}-v_{\text{медл}}}.
\]

Если между объектами есть начальный промежуток $d$,
его добавляют к относительному расстоянию:

\[
S_{\text{отн}}=d+L_1+L_2.
\]

\subsect{Единицы}

Если длины заданы в метрах, можно перевести скорость в м/с:

\[
1\unit{км/ч}=\frac{5}{18}\unit{м/с}.
\]

Либо можно перевести метры в километры и оставить скорость в км/ч.
Оба способа равноправны при согласованных единицах.

% ============================================================
% ПРОИЗВОДИТЕЛЬНОСТЬ
% ============================================================

\sect{8. Производительность}

Для задач на работу действует та же структура:

\[
A=pt,
\qquad
t=\frac{A}{p},
\qquad
p=\frac{A}{t}.
\]

Если два исполнителя выполняют разные объёмы работы за одинаковое время,
то

\[
\frac{A_1}{p_1}=\frac{A_2}{p_2}.
\]

\subsect{Пример}

Первая труба наполняет бак объёмом $770\unit{л}$,
а вторая --- бак объёмом $830\unit{л}$ за то же время.
Одна из труб пропускает на $6\unit{л/мин}$ больше другой.

Поскольку за одинаковое время вторая труба наполняет больший бак,
её производительность выше.
Пусть производительность первой трубы равна $x\unit{л/мин}$.
Тогда производительность второй --- $(x+6)\unit{л/мин}$, причём $x>0$.

Из равенства времён:

\[
\frac{770}{x}=\frac{830}{x+6}.
\]

Знаменатели положительны, поэтому можно перемножить крест-накрест:

\[
770(x+6)=830x,
\]

\[
770x+4620=830x,
\]

\[
4620=60x,
\qquad
x=77.
\]

Следовательно, первая труба пропускает

\[
77\unit{л/мин},
\]

а вторая ---

\[
83\unit{л/мин}.
\]

% ============================================================
% ТИПИЧНЫЕ ОШИБКИ
% ============================================================

\sect{9. Типичные ошибки}

\begin{itemize}
    \item Начинать вычисления до того, как определено, что означает $x$.
    \item Путать сумму и разность скоростей: навстречу используется сумма, вдогонку --- разность.
    \item Игнорировать более ранний старт одного участника.
    \item На круговой трассе забывать условие обгона: $S_{\text{быстр}}-S_{\text{медл}}=L$.
    \item При движении плота использовать собственную скорость лодки вместо скорости течения.
    \item Находить среднюю скорость как простое среднее скоростей без проверки структуры пути.
    \item В задачах на протяжённые тела забывать длину второго объекта или начальный промежуток.
    \item Смешивать километры и метры, часы и минуты внутри одного уравнения.
\end{itemize}

% ============================================================
% ТРЕНИРОВОЧНЫЙ БЛОК
% ============================================================

\clearpage

\sect{Тренировочный блок · Путь А}

Решите задачи табличным методом. Для каждой задачи:
обозначьте неизвестную, составьте таблицу, запишите уравнение,
решите его, проверьте ограничения и только затем сформулируйте ответ.

\begin{enumerate}

\item
Из городов $A$ и $B$, расстояние между которыми $360\unit{км}$,
одновременно навстречу друг другу выехали два автомобиля
со скоростями $60\unit{км/ч}$ и $90\unit{км/ч}$.
На каком расстоянии от города $A$ они встретятся,
если первый автомобиль выехал из $A$?

\item
Расстояние между городами равно $450\unit{км}$.
Первый автомобиль выехал из города $A$ со скоростью $60\unit{км/ч}$
на $30\unit{мин}$ раньше второго,
который выехал из города $B$ навстречу ему со скоростью $90\unit{км/ч}$.
На каком расстоянии от города $A$ произойдёт встреча?

\item
Медленный пешеход идёт со скоростью $5\unit{км/ч}$.
Быстрый пешеход выходит вслед за ним из той же точки,
когда первый уже находится на расстоянии $450\unit{м}$,
и движется со скоростью $6{,}5\unit{км/ч}$.
Через сколько минут быстрый пешеход догонит медленного?

\item
По круговой трассе длиной $12\unit{км}$ одновременно
из одной точки в одном направлении выехали два автомобиля.
Скорость быстрого автомобиля равна $78\unit{км/ч}$.
Через $40\unit{мин}$ он впервые обогнал медленного ровно на один круг.
Найдите скорость медленного автомобиля.

\item
Два велосипедиста стартовали из противоположных точек
круговой трассы длиной $20\unit{км}$
и поехали навстречу друг другу со скоростями
$12\unit{км/ч}$ и $8\unit{км/ч}$.
Через сколько минут они впервые встретятся?

\item
Моторная лодка проходит некоторое расстояние по течению
со скоростью $30\unit{км/ч}$, затем стоит $2\unit{ч}$
и возвращается против течения со скоростью $20\unit{км/ч}$.
Весь рейс занимает $7\unit{ч}$.
Найдите суммарный путь лодки туда и обратно.

\item
Плот движется по реке со скоростью течения $3\unit{км/ч}$.
Через $2\unit{ч}$ после отправления плота из той же точки
вслед за ним вышла моторная лодка.
Её собственная скорость равна $15\unit{км/ч}$.
На каком расстоянии от пункта отправления лодка догонит плот?

\item
Первую треть пути автомобиль ехал со скоростью $12\unit{км/ч}$,
вторую треть --- со скоростью $18\unit{км/ч}$,
последнюю треть --- со скоростью $36\unit{км/ч}$.
Найдите среднюю скорость автомобиля на всём пути.

\item
Поезд длиной $180\unit{м}$ движется со скоростью $72\unit{км/ч}$
и полностью проходит мимо платформы длиной $270\unit{м}$.
Сколько секунд длится прохождение платформы
от момента, когда локомотив поравнялся с началом платформы,
до момента, когда последний вагон миновал её конец?

\item
Первая труба наполняет бак объёмом $720\unit{л}$ за то же время,
за которое вторая труба наполняет бак объёмом $900\unit{л}$.
Вторая труба пропускает на $10\unit{л/мин}$ больше первой.
Сколько литров воды в минуту пропускает первая труба?

\end{enumerate}

% ============================================================
% ОТВЕТЫ
% ============================================================

\clearpage

\sect{Ответы}

\begin{table}[h!]
\centering
\begin{tabularx}{0.9\linewidth}{@{}cX@{}}
\toprule
\textbf{№} & \textbf{Ответ} \\
\midrule
1 & $144\unit{км}$ \\
2 & $198\unit{км}$ \\
3 & $18\unit{мин}$ \\
4 & $60\unit{км/ч}$ \\
5 & $30\unit{мин}$ \\
6 & $120\unit{км}$ \\
7 & $7{,}2\unit{км}$ \\
8 & $18\unit{км/ч}$ \\
9 & $22{,}5\unit{с}$ \\
10 & $40\unit{л/мин}$ \\
\bottomrule
\end{tabularx}
\end{table}

\vfill

\begin{center}
\small\color{softgray}
Перед сверкой ответа проверьте смысл составленного уравнения,
ограничения и единицы измерения.
\end{center}

\end{document}